\begin{abstract}
未来空域需求感知需要在较长提前量上估计航空器可能占用的区域，并给出可用于聚合、排序和风险分析的概率分布。
本文研究多时域层级区域航空轨迹概率预测：基于 ADRR 派生的一个月商业航班轨迹数据构建 20,712,602 个预测窗口，给定按 5 分钟重采样的 30 分钟历史轨迹，同时预测未来 10、30、60 和 120 分钟的 H3 区域概率。
为避免把长时不确定性压缩成单一路径，本文采用 coarse-to-fine H3 输出协议，先预测父层区域，再在父层条件下恢复细粒度区域概率。
方法上，本文提出多尺度 horizon-query 编码器，使每个预测时域分别从细尺度历史序列和 stride-pooled 粗尺度序列中检索时间证据。
在此基础上，本文提出 Horizon-Router 单模型，用 global、coarse-stable 和 local-causal 三个内部 expert 表达不同时间尺度的运动证据，并由 horizon-level router 为每个预测时域自适应融合。
训练阶段进一步引入基于 expert loss 的 router supervision、活跃 expert mask、负载均衡、expert auxiliary loss 和 horizon-wise teacher distillation，以缓解专家塌缩并稳定大规模训练。
完整数据集结果显示，Horizon-Router 的宏平均 top-1/top-5 达到 0.6884/0.9710，NLL、Brier 和 ECE 分别为 0.9399、0.4262 和 0.0104；相较适配后的 TrAISformer-H3 强基线，模型在 top-1、top-5、NLL 和 Brier 上均取得提升。
这些结果表明，面向未来空域需求感知的长时航空预测可以被有效表述为层级区域概率建模问题，而时域级专家路由为不同提前量下的历史证据选择提供了有效结构。
\end{abstract}

\noindent\textbf{关键词：} 航空轨迹预测；空域需求感知；H3 空间索引；多时域预测；概率预测；专家路由
